Our findings reveal both the promise and the current limitations of data-driven models for age-related motion signatures.
The ST-GCN successfully separates the kinematic extremes (Young vs.\ Elderly), confirming that motion clips contain discriminative biomechanical information related to age group.
The LoRA-adapted MDM learns and imposes spatial kinematic constraints that the unadapted MDM cannot express even under elderly-described prompts, yet spatiotemporal consistency remains unresolved: generated stride variability massively overshoots clinical baselines.
We attribute this remaining failure to the capacity ceiling of low-rank, discrete-token adaptation: a rank-12 LoRA update can push high-leverage spatial statistics such as hip RoM toward an age-specific mean, but cannot coherently re-shape stride timing, leaving the base MDM prior to dominate the temporal axis of the kinematic chain.
This gap highlights a core challenge for assistive robotics in eldercare: visually plausible age-conditioned motion is not sufficient --- prediction systems built on top of generative priors must preserve the biomechanical characteristics that affect anticipation, safety, and interaction.

Dataset scale and discrete-token conditioning capacity are the two principal bottlenecks; future work will train the continuous latent-conditioning architecture of Section~\ref{sec:method} end-to-end on larger, age-diverse datasets, with scalar age regression and t-SNE visualization of $\mathbf{z}_\text{age}$ as classifier-side next steps --- moving past the LoRA capacity ceiling toward demographic-aware motion priors for assistive eldercare.